\documentclass[12pt,a4paper]{ctexart}
\usepackage{amsmath,amssymb}
\usepackage{booktabs}
\usepackage{array}
\usepackage{geometry}
\usepackage{enumitem}
\usepackage[hidelinks]{hyperref}
\geometry{margin=2.4cm}
\setlength{\parskip}{0.4em}
\setlength{\parindent}{2em}
\linespread{1.25}
\renewcommand{\arraystretch}{1.2}
\setcounter{secnumdepth}{-1}
\ctexset{
  section = {format = {\heiti \zihao{3}}},
  subsection = {format = {\heiti \zihao{4}}},
}

\begin{document}

\begin{center}
{\heiti\zihao{-2} External 新增资料可用性评估报告\par}
\vspace{0.4em}
{\heiti\zihao{4} ——四份文献与论文《A题论文》的对接点汇总（供评审，未修改论文）\par}
\end{center}

\section{一、评估原则}

按既定定位：本题是\textbf{建模与计算}任务，模型输入只使用题给附录经验式与附件数据；外部资料中的实测数据\textbf{只用于验证对照，不进入模型、不替换题给参数、不改变任何数值结果}。据此，四份资料的用途被分为三类：

\begin{itemize}[leftmargin=2em]
\item \textbf{A 类·理论依据}：标准教材/文献中的定理、公式、机理表述，用于支撑论文的分析与检验段落（如湿球温度方程、干燥两阶段机理），以引用方式吸收；
\item \textbf{B 类·验证对照}：文献实测数据与本文计算结果的交叉比对（如 $D_{\mathrm{eff}}$ 量级、$E_a$ 区间），只出现在"模型检验"一节，不进入模型；
\item \textbf{C 类·术语规范}：用于校核论文措辞是否与学科标准用法一致（本次发现一处需要复核，见 3.2 节）。
\end{itemize}

\section{二、资料总览}

\begin{center}
{\zihao{5}\begin{tabular}{llll}
\toprule
文件 & 类型 & 规模与可提取性 & 总体定位 \\
\midrule
任迪峰博士学位论文（2004，中国农业大学） & 学位论文 & 105 页，有 OCR 文字层（数字多处乱码） & B 类为主，兼 A/C 类 \\
陈敏恒《化工原理》（第五版）下册 & 教材 & 277 页，文字层完整 & A 类（湿球温度、固体干燥） \\
陈敏恒《化工原理》（第五版）上册 & 教材 & 321 页，文字层完整 & A 类（对流传热关联式） \\
陶文铨《传热学》（第五版） & 教材 & 591 页，扫描版无文字层，仅书签目录可用 & A 类（仅可作引用出处） \\
\bottomrule
\end{tabular}}
\end{center}

\section{三、分册评估}

\subsection{3.1 任迪峰《中药材干燥过程中质量退化及优化干燥工艺的研究》（已为论文文献 [9]，本次获得全文）}

\textbf{内容概要}：以地黄、麦冬为对象，测定了解吸平衡含水率（四种平衡含水率模型拟合）、对流换热系数与当量导热系数（由薄层干燥试验能量衡算反算）、水分扩散系数（非稳态扩散方程反算）、收缩系数与孔隙率、比热；建立含\textbf{导湿温性}（雷科夫/Luikov 型）项的传热传质微分方程，用有限元法求解；完成变风温、发汗干燥工艺多目标优化。

\textbf{可用点（按论文位置）}：

\begin{enumerate}[leftmargin=2em]
\item \textbf{[B 类，论文 7.5]} 实测 $D_{\mathrm{eff}}$ 的定性规律：温度越高、含水率越高，$D_{\mathrm{eff}}$ 越大（论文第 3.4 节，PDF 第 48--49 页）。这与题给 $D(C,T)$ 随 $C$ 增大的定性趋势一致，可作为 7.5 节"参数物理合理性"的补充互证。具体数值因 OCR 乱码需人工核对原书页（书页码：3.2 节 p.40、3.3 节 p.44、3.4 节 p.47、3.5 节 p.49、3.6 节 p.52；PDF 页 $=$ 书页 $+1$），\textbf{如需引用具体数值必须人工核对}。
\item \textbf{[A 类，论文 7.6/H5]} 该论文反算对流换热系数 $h$ 的能量衡算中\textbf{显式包含汽化潜热项}（p.44），即文献实践中的 $h$ 是从含潜热的能量平衡里回归出的\textbf{有效值}——与本文假设 H5"经验参数已吸收汽化效应"的观点一致，可作为 H5 的文献先例佐证。
\item \textbf{[A 类，论文 8.2(1)]} 其水分扩散方程含\textbf{导湿温性项}（温度梯度驱动的水分迁移，雷科夫理论，p.57--58），正是本文 8.2(1) 所述"Luikov 型模型"的具体实例：引用时应表述为\textbf{扩展方向}（本文模型按题设无 Soret/Dufour 交叉项，二者不可混同）。
\item \textbf{[A 类，论文 8.3]} 变温干燥机理结论（p.101）：干燥后期\textbf{降温}使物料内部湿度梯度与温度梯度方向一致，促使水分进一步脱出、出现恒速段，且平均料温较低、有利于防止有效成分降解；发汗过程中温度越高、有效成分含量下降越快。这直接支撑本文 8.3 节"变温策略（分段降温保成分）优化"与"高温高湿区降解风险"的表述，建议在该段补引 [9]。
\item \textbf{[C 类，论文 6.3]（重要）} 该论文对地黄薄层干燥的实测结论：\textbf{没有出现恒速干燥阶段}，干燥主要发生在降速段（p.101）。对照教材定义（恒速段要求表面为非结合水覆盖、表面温度稳定在湿球温度，见 3.2 节），本文表 3/表 4 中表面温度在 0.5--1.5 h 内持续上升、表面含水率持续下降，严格说不构成恒速段。\textbf{建议将 6.3 节"表面传质主导的恒速干燥段"改为"表面传质主导的快速干燥段"或加注"近似恒速、无典型恒速段"}，使术语与学科标准一致。
\item \textbf{[A 类，论文 H2]} 扩散系数试验中将圆柱体两端涂环氧树脂以消除端面效应、保证径向一维扩散（p.48）——与本文假设 H2（径向一维）的实验做法一致，可作 H2 的文献佐证。
\item \textbf{[A 类，论文 6.3 平衡含水率]} 解吸平衡含水率模型（GAB 等四模型比较，p.41--43）可在 6.3 节"趋近平衡含水率"处补引，低优先级。
\end{enumerate}

\subsection{3.2 陈敏恒《化工原理》下册（第五版）：第 13 章热质同时传递、第 14 章固体干燥}

\textbf{可用点}：

\begin{enumerate}[leftmargin=2em]
\item \textbf{[A 类，论文 7.6]（重要）} 湿球温度的标准理论：教材式 (13-7) 即"显热$=$潜热"平衡
\begin{equation}
\alpha(t-t_w)=k_H(H_w-H)\,r_w,\qquad
t_w=t-\frac{k_H}{\alpha}r_w(H_w-H),
\end{equation}
并给出空气--水系统 $\alpha/k_H\approx1.09$ kJ/(kg$\cdot$$^\circ$C)（流速 $>5$ m/s）。论文 7.6 节目前写"由焓湿图（或湿球温度方程）读得 $T_{wb}\approx40\ ^\circ$C"，建议补引教材 13.2.2 节，并可加一句核算：\textbf{按式 (13-9) 迭代（$p_s$ 查水蒸气表）得 $t_w\approx41$--$42\ ^\circ$C}，与论文取值 40 $^\circ$C 相差约 1.5 $^\circ$C，对 7.6 节敏感性结论（$D$ 下降约 32\%、时长约 $+30\%$）无实质影响；若求严谨可写"约 40--42 $^\circ$C"。此核算属 B 类验证，不改变模型与结果。
\item \textbf{[A 类，论文 6.3]（术语依据）} 干燥两阶段的标准定义与机理：恒速段（表面为非结合水覆盖，\textbf{表面温度即湿球温度且维持不变}）、降速段四原因（实际汽化表面减小、汽化面内移、平衡蒸气压下降、内部扩散极慢）、临界含水量定义与典型值（表 14-1）。这是论文 6.3 节两段干燥表述的标准出处，建议引用第 14.3.1 节。
\item \textbf{[A 类，论文 8.3]} 工艺原则："在恒速阶段，物料表面温度维持在湿球温度，因此即使在高温下易于变质、破坏的物料（塑料、药物、食品等）仍然允许在恒速阶段采用较高的气流温度……在降速阶段，物料温度逐渐升高，故在干燥后期须注意不使物料温度过高"——直接支撑本文 8.3 节的变温/分段控温推广表述。
\item \textbf{[A 类，论文 8.2/8.3]} 间歇干燥时间估算：恒速段时间 $\tau_1=(G_c/A)(X_1-X_c)/N_A$ 及降速段计算（14.3.2 节）——可作为"工程简化算法"与本文扩散模型对照的推广/评价素材（论文 8.3 或评价节可选）。
\item \textbf{[A 类，论文 H4]} 恒速段传质速率式 $N_A=k_H(H_w-H)$（14-11）中 $k_H$ 在定态空气条件下为常数——与本文假设 H4（$h,h_m$ 取常数）的物理基础一致，可补一句教材依据。
\end{enumerate}

\textbf{引用格式建议}：陈敏恒, 丛德滋, 方图南, 等. 化工原理（下册）[M]. 5 版. 北京: 化学工业出版社, 2020.

\subsection{3.3 陈敏恒《化工原理》上册（第五版）：第 6 章传热}

\begin{enumerate}[leftmargin=2em]
\item \textbf{[A 类，论文 8.2(3)]} 6.3.3 节"无相变的对流给热系数的经验关联式"：给出强制对流的 $Nu=f(Re,Pr)$ 关联式体系。论文 8.2(3) 提出"可用经验关联式（如 Chilton--Colburn 类比）修正 $h,h_m$"，可补引上册 6.3.3 节作为具体出处。
\item \textbf{[A 类，论文 2.2/教程]} 6.2.1 节傅里叶定律与 6.2.3 节圆筒壁稳态导热——与论文第二章推导同源，可作推导的标准教材依据（优先级低，传热学 [13] 已覆盖）。
\item \textbf{[C 类，术语]} 6.6.5 节"非定态传热过程的拟定态处理"——若论文将来讨论"拟定态假设"需区分此概念（当前论文未使用，无需改动）。
\end{enumerate}

\textbf{引用格式建议}：陈敏恒, 丛德滋, 方图南, 等. 化工原理（上册）[M]. 5 版. 北京: 化学工业出版社, 2020.

\subsection{3.4 陶文铨《传热学》第五版（扫描版，无文字层）}

书签目录显示各章起页：第 2 章稳态热传导（p.51）、\textbf{第 3 章非稳态热传导（p.119，含无限长圆柱的级数解析解——Bessel 解的教材出处）}、第 4 章热传导问题的数值解法（p.171，差分格式与追赶法）、第 5--6 章对流传热（p.207/253，$h$ 的准则关联式）、第 7 章相变对流传热（p.313）、第 11 章传质学简介（p.537）。

\begin{enumerate}[leftmargin=2em]
\item \textbf{[A 类，论文 7.1]} 论文 Bessel 解析解目前引 Carslaw--Jaeger 专著 [2]；建议\textbf{补充或更新}教材引用：陶文铨. 传热学[M]. 5 版. 北京: 高等教育出版社, 2019.（第 3 章无限长圆柱非稳态导热级数解），使解析解对照有标准教材出处。论文现有 [13] 为第四版（2006），可更新为第五版并具体到章。
\item \textbf{[A 类，教程附录 C]} 第 4 章数值解法与教程《A题方法详解》第 5--7 章同源，可作为教程学习路线的教材条目（与陶文铨《数值传热学》并列）。
\item 因无文字层，\textbf{无法提取具体公式与页码核对}，引用时章节号依据书签目录，具体页码需对照纸质书。
\end{enumerate}

\section{四、汇总建议表}

\begin{center}
{\zihao{5}\begin{tabular}{p{4.6cm}p{1.9cm}p{3.2cm}p{3.6cm}p{1.3cm}}
\toprule
建议内容 & 类别 & 对应论文位置 & 建议动作 & 优先级 \\
\midrule
6.3 节"恒速干燥段"措辞 & C & 6.3 节 & 改为"表面传质主导的快速干燥段"或加"近似恒速"注 & 高 \\
湿球温度方程补引与核算 & A/B & 7.6 节 & 补引化工原理下册 13.2.2 节式 (13-9)；可加"迭代核算约 41--42 $^\circ$C" & 高 \\
恒速段"表面温度=湿球温度"、降速四原因、临界含水量 & A & 6.3 节 & 补引化工原理下册 14.3.1 节 & 高 \\
"恒速段允许高温、降速段控温"工艺原则 & A & 8.3 节 & 补引化工原理下册 14.3.1 节 + 任迪峰变温结论 [9] & 中 \\
Luikov 型（导湿温性）扩展方向具体化 & A & 8.2(1) 节 & 引用任迪峰模型（p.57--58）为实例，注明为扩展方向 & 中 \\
$h$ 反算含潜热项 = 有效参数先例 & A & 7.6/H5 & 补引任迪峰 p.44 做法佐证 H5 & 中 \\
$D_{\mathrm{eff}}$ 随 $C,T$ 变化趋势互证 & B & 7.5 节 & 选择性补充（数值需人工核对 OCR） & 低 \\
Bessel 解教材出处 & A & 7.1 节 & [13] 更新为传热学第五版并引第 3 章 & 低 \\
对流给热系数关联式出处 & A & 8.2(3) 节 & 补引化工原理上册 6.3.3 节 & 低 \\
端面涂胶保证径向一维 & A & H2 & 可选补引任迪峰 p.48 & 低 \\
平衡含水率模型（GAB 等） & A & 6.3 节 & 可选补引任迪峰 3.2 节 & 低 \\
工程干燥时间估算法（$\tau_1$ 公式） & A & 8.3/评价节 & 可选提及作为对照 & 低 \\
\bottomrule
\end{tabular}}
\end{center}

\section{五、结论}

\textbf{值得吸收的}（均为"引用与措辞"层面，不改模型、不改参数、不改任何数值结果）：

\begin{enumerate}[leftmargin=2em]
\item 论文 6.3 节"恒速干燥段"的措辞应复核——教材严格定义与任迪峰实测均不支持典型恒速段，建议改为"表面传质主导的快速干燥段"（术语规范收益最大的一处）；
\item 论文 7.6 节湿球温度分析可补引化工原理下册 13.2.2 节的严格推导，并可用式 (13-9) 迭代核算（得约 41--42 $^\circ$C，与现取 40 $^\circ$C 相差约 1.5 $^\circ$C，敏感性结论不变；核算仅作验证，不进模型）；
\item 论文 6.3/8.3 节的两段干燥机理与变温工艺论述可补引化工原理下册第 14 章与任迪峰变温干燥结论 [9]，使表述有标准出处；
\item 论文 8.2(1) 节的 Luikov 改进方向可用任迪峰模型（含导湿温性项）具体化；
\item 参考文献可新增化工原理上、下册（第五版）两条，并将传热学 [13] 更新为第五版、具体到章。

\textbf{不动的}：模型结构、题给附录参数、附件数据处理、全部表格与 $t^*$ 数值——四份资料中没有任何内容要求或支持改变这些；任迪峰的实测物性数据仅用于 7.5 节的文献互证（且引用具体数值前必须人工核对 OCR 原页）。

\textbf{附注（提取质量说明）}：任迪峰论文为扫描 OCR，正文文字可读但\textbf{数字多处乱码}，本报告引用其定性结论与公式结构，不引用具体数值；陶文铨《传热学》第五版无文字层，仅书签目录可用；化工原理两册为原生文字版，公式与数据可靠。本文未修改《A题论文.tex》，以上全部为建议，供逐条审定后另行执行。

\end{document}
